\section{Introduction}\label{sec:introduction}

An insertion or deletion changes the positions of every later symbol.
Yet cyclically shifting a nonempty word by one position costs at most two edits: delete its first
symbol and insert it at the end. We use the difference between the length of
a word and the cost of translating it to study how faithfully edit distance
can be represented in $\ell_1$.

\paragraph{The metric and the question.}
For finite words $x,y$, let $\ED(x,y)$ be the least number of single-symbol
insertions, deletions, and substitutions taking $x$ to $y$, with each
operation costing one. If $\Sigma$ is a finite alphabet with at least two
symbols and $d\geq1$ is an
integer, let $\Sigma^{\leq d}$ contain all words of length at most $d$,
including the empty word. This cap restricts only the endpoints of a
comparison; intermediate words in an edit script are unrestricted.
For an injective map $f:\Sigma^{\leq d}\to\ell_1$, its distortion is
\[
 \operatorname{dist}(f)=
 \left(\max_{x\ne y}\frac{\|f(x)-f(y)\|_1}{\ED(x,y)}\right)
 \left(\max_{x\ne y}\frac{\ED(x,y)}{\|f(x)-f(y)\|_1}\right),
 \qquad E_\Sigma(d)=\inf_f\operatorname{dist}(f).
\]
Here $\ell_1$ is the real space of absolutely summable sequences. There is
no restriction on target dimension or on how the map is computed. The
question asks for one map that represents every pair in the domain.

\begin{theorem}\label{thm:main}
There are absolute constants $c,C>0$ and an integer $d_0$ such that, for
every integer $d\geq d_0$ and every finite alphabet $\Sigma$ with
$|\Sigma|\geq2$,
\begin{equation}\label{eq:main}
 \exp\!\bigl(c\sqrt{\log d\,\log\log d}\bigr)
 \leq E_\Sigma(d)\leq
 \exp\!\bigl(C\sqrt{\log d\,\log\log d}\bigr).
\end{equation}
Each of our two lower constructions supplies a binary subset whose words
have one common length at most $d$ and whose least $\ell_1$ distortion
satisfies the lower bound. The upper bound is realized by a deterministic
finite-dimensional map.
\end{theorem}

All logarithms are natural unless a base is displayed. The constants are
independent of the finite alphabet, including alphabets whose size grows
with $d$. Thus the same estimates hold for
$\sup_{2\leq|\Sigma|<\infty}E_\Sigma(d)$. Equation~\eqref{eq:main}
determines the order of $\log E_\Sigma(d)$; the constants in its two
exponents need not agree.

\paragraph{Earlier bounds and the role of alignment.}
The single-symbol error model has its roots in Levenshtein's work on
insertion and deletion codes \cite{Levenshtein1966}. Wagner and Fischer's
increasing traces describe the alignments available to two words
\cite{WagnerFischer1974}. In particular, if $\LCS(x,y)$ denotes the longest
common subsequence length, the insertion--deletion distance is
\[
 \Delta(x,y)=|x|+|y|-2\LCS(x,y),\qquad
 \ED(x,y)\leq\Delta(x,y)\leq2\ED(x,y).
\]
A lower bound must therefore control arbitrary increasing matchings, even
when they cross many of the boundaries used to construct the words.

Ostrovsky and Rabani obtained the upper scale in Theorem~\ref{thm:main} for
fixed-length binary edit distance by recursively comparing multisets of
overlapping substrings \cite{OstrovskyRabani2007}. They also noted extensions
to larger alphabets and varying lengths. Section~\ref{upper:section} gives
a complete finite histogram construction with alphabet-independent constants
and all shorter words included. It can be read on its own.

For lower bounds, Andoni, Deza, Gupta, Indyk, and Raskhodnikova found binary
sets whose $L_1$ distortion approaches $3/2$
\cite{AndoniDezaGuptaIndykRaskhodnikova2003}. Their construction uses several
nearby word lengths. Khot and Naor obtained a
$(\log d)^{1/2-o(1)}$ bound using Fourier analysis of cuts and noisy shifts
\cite{KhotNaor2006}; their insertion--deletion convention transfers to ours
within a factor of two. Krauthgamer and Rabani strengthened the bound to
$\Omega(\log d)$ on binary words of length $d$
\cite[Theorem~1.3]{KrauthgamerRabani2009}. Our constructions use the same
broad cut and Fourier viewpoint, but impose translation constraints through
distinct prime periods. Those constraints force large total leaf frequency
and yield the exponent in~\eqref{eq:main}. The resulting lower bound grows
faster than every fixed power of $\log d$, while remaining $d^{o(1)}$.
Thus it reaches the scale of the Ostrovsky--Rabani upper bound after taking
logarithms.

\paragraph{The lower constructions.}
A phase is a point of the circle $\mathbb R/\mathbb Z$. We work on a finite
equally spaced grid containing the half-turn $1/2$. A state assigns one grid
phase to each leaf of a finite rooted tree. At a leaf, a binary pattern
records a translated half-circle on that grid. At an internal node, a row
lists the child words, each with its own tag; a long word repeats rows at
translated phases.

Both constructions compare three kinds of phase change. Translations in a
specified subgroup move the row lists at low edit cost. Moving just one
leaf changes only the fraction of the word occupied by that leaf, with
edit cost at most twice the occupied length times the phase displacement.
A simultaneous half-turn at every leaf must have a large edit cost. This last
assertion requires control of arbitrary alignments: a matching may cross
many of the rows used to define the words.

The row tags provide that control. A source row whose matches stay in one
target row inherits the child separation. If its matches enter several
target rows, they leave a long unmatched target interval. The intervals
charged by distinct source rows are disjoint. Lemma~\ref{pre:tagged-rows}
turns this observation into the separation estimate used by both
constructions.

To compare these edit costs with an arbitrary $\ell_1$ image, we express
its distances on the finite phase space as nonnegative sums of cut distances
and then use finite Fourier expansion. A character has one integer frequency
at each leaf. Averaging over the cheap subgroup controls the total Fourier
weight of characters nonconstant on it. Among the remaining characters,
one that detects the half-turn must have large total absolute leaf
frequency: if that total were small, the prime constraints would force
nonzero frequency sums to branch through too many leaves. Weighted
single-leaf shifts control the Fourier mass carried by these large
frequencies.
The common analytic step is Lemma~\ref{pre:weighted-circle}; each construction proves
its own arithmetic branching bound.

The two prime patterns also give different proofs of half-turn separation.
The first construction groups children by primes and includes every residue
of each prime. Its separation proof follows a nonzero phase down the tree
while keeping track of the loss at each level. The complete residue lists
allow the descent to restart at the reciprocal of a single prime; these
resets keep the denominator factors acquired at other steps from
accumulating without control. The accumulated losses and the final phase's
distance from zero give a quantitative lower bound for the half-turn.
Section~\ref{sec:masks} constructs the words and carries out this descent.

The second construction uses one child per prime, with that prime's
translation acting on every child except its own. Two child phases of small
cost would have a small difference involving at most two fresh prime
denominators. The induction rules out every nonzero such difference, forcing
most translation coordinates to vanish. Section~\ref{sec:excluded} develops
this separate argument.

\paragraph{Binary transfer and the upper construction.}
The tags initially require a larger alphabet. Section~\ref{p2:sec:binary-coding}
gives two direct binary substitutions of width $O(\log(2A))$ for an alphabet
of size $A$. Each has absolute distortion on all equal-length words,
simultaneously for every common length. One proof separates intervals of
differently labelled codewords and uses a common offset to recover a letter
matching. The other controls shifted self-overlaps and finds anchors at equal
relative indices. The lower constructions use these bounds only at equal lengths.

The same existence order follows from the alphabet reduction of
Bhattacharya, Dey, Goldenberg, and Kouck\'y \cite[Theorem~1.1 and Section~3.2.1 of the full version]{BhattacharyaDeyGoldenbergKoucky2024}.
Their stated theorem approximately preserves normalized insertion--deletion
distance, with each distance divided by the sum of the lengths of its two
arguments. The block substitution in its proof has logarithmic width for a
fixed target alphabet. Composing that substitution
with a fixed-width binary delimiter gives the claimed binary width and
distortion, as we explain in Section~\ref{p2:sec:binary-coding}. Our direct proofs expose two different
local matching conditions. Long-interval separation is already present
in Schulman and Zuckerman's inner code \cite[Section~II-C]{SchulmanZuckerman1999},
and self-matching control is a central feature of synchronization strings
\cite{HaeuplerShahrasbi2017}.

For either lower route, depth $h$ gives logarithmic distortion of order
$h\log h$ at lengths whose logarithms are $O(h^2\log h)$. The direct
parameter choice in Section~\ref{sec:assembly} supplies every sufficiently
large cap. Finally, Section~\ref{upper:section} develops the independent
upper construction: divide words into blocks, list overlapping windows,
partition their shorter-word images using cuts, and record the resulting
histograms. After the power-length proof, a map at actual window lengths
admits two expansion estimates, one from an arbitrary longest common
subsequence and one from moves. A final variant uses total variation.

Theorem~\ref{thm:uniform-upper} supplies the upper embedding stated as
Theorem~6.1 of the companion tree constructions \cite{OpenAITree2026},
which is combined with their lower bounds in Corollary~6.2. Their lower
proofs are independent of the present constructions.
A shorter companion proof \cite{OpenAIMain2026} uses marked rows and a
prime-product displacement inequality, with no lower-proof dependence on
the circle constructions.
